\chapter{Introdução}\label{chap:introducao}
Este documento apresenta um exemplo de dissertação, tese ou relatório de qualificação do PPGMC/UESC. A versão 2.0 preserva a vocação didática da versão 1.4: além de servir como ponto de partida, demonstra os principais recursos de escrita científica em \LaTeX. O conteúdo exemplificativo deve ser substituído pelo trabalho do estudante.

A modelagem computacional permite investigar fenômenos por meio de modelos matemáticos, discretizações e algoritmos. Um problema científico envolve perguntas de pesquisa, hipóteses, dados, métodos e critérios de validação. A introdução deve delimitar o problema, explicitar sua relevância e apresentar o percurso argumentativo.

\section{Contextualização e motivação}
A descrição do problema deve indicar o contexto científico, o estado do conhecimento e a lacuna investigada. Recomenda-se fundamentar afirmações relevantes com fontes bibliográficas verificáveis. A construção de modelos deve explicitar hipóteses, escalas e limites de aplicabilidade.

\section{Problema de pesquisa e justificativa}
O estudante deve formular uma pergunta clara, justificar sua relevância e identificar as contribuições esperadas. Evite confundir a utilização de uma ferramenta com o objetivo científico da investigação.

\section{Objetivos}
\subsection{Objetivo geral}
Investigar, por meio de modelagem computacional, um fenômeno de interesse, explicitando hipóteses e critérios de avaliação.
\subsection{Objetivos específicos}
\begin{enumerate}[label=\alph*)]
\item Formular um modelo matemático apropriado ao problema.
\item Implementar e verificar uma solução computacional.
\item Avaliar os resultados e as limitações do método.
\end{enumerate}

\section{Interdisciplinaridade da pesquisa}
A interdisciplinaridade deve ser apresentada como integração efetiva de contribuições, e não apenas como uma lista de áreas. No exemplo da Figura~\ref{fig:interdisciplinaridade}, a Física descreve fenômenos, a Matemática fornece modelos e métodos, a Computação viabiliza algoritmos e a Engenharia Nuclear estabelece aplicações e requisitos de validação. O estudante deve substituir as áreas e as contribuições pelas que realmente integram seu trabalho.
\begin{figure}[htbp]\centering
\resizebox{\textwidth}{!}{% Edite os quatro pares: nome da área / contribuição.
\begin{tikzpicture}[>=Latex, every node/.style={align=center,font=\small},
 area/.style={draw,rounded corners=9pt,minimum width=3.7cm,minimum height=1.3cm,thick},
 central/.style={circle,draw=black!65,fill=blue!13,minimum size=3cm,thick}]
\node[central] (p) at (0,0) {\textbf{Esta pesquisa}};
\node[area,draw=blue!65,fill=blue!7] (a) at (-3.65,2.35) {\textbf{Física}\\Fenômenos e leis};
\node[area,draw=green!55!black,fill=green!7] (b) at (3.65,2.35) {\textbf{Matemática}\\Modelos e métodos};
\node[area,draw=orange!80!black,fill=orange!7] (c) at (-3.65,-2.35) {\textbf{Computação}\\Algoritmos e simulação};
\node[area,draw=violet,fill=violet!7] (d) at (3.65,-2.35) {\textbf{Engenharia Nuclear}\\Aplicações e validação};
\draw[->,very thick,blue!65] (a) -- (p);
\draw[->,very thick,green!55!black] (b) -- (p);
\draw[->,very thick,orange!80!black] (c) -- (p);
\draw[->,very thick,violet] (d) -- (p);
\end{tikzpicture}
}
\caption{Exemplo editável de integração interdisciplinar em uma pesquisa.}\label{fig:interdisciplinaridade}
\fonte{Elaboração própria (exemplo didático).}
\end{figure}

\section{Organização do documento}
O Capítulo~\ref{chap:organizacao} explica os arquivos do projeto; o Capítulo~\ref{chap:citacoes} apresenta citações e referências; o Capítulo~\ref{chap:elementos} reúne figuras, quadros, tabelas, equações e algoritmos; os capítulos seguintes exemplificam metodologia, resultados e conclusão. No modo qualificação, acrescenta-se um cronograma de atividades.
